\chapter{Eulero}

Una volta realizzata la mesh in base alle caratteristiche del problema in esame, come descritto nel capitolo precedente, si sfrutta il risultato del software e lo si immette all'interno del codice CFD. Il codice che si intende sviluppare risolve le equazioni di Eulero 2D e il linguaggio di programmazione utilizzato è Fortran 90 date le sue alte prestazioni.

\section{Strutture}

Il codice, per funzionare, ha bisogno di salvare i dati che arrivano dal file di gmsh in alcune variabili, che andranno opportunamente definite all'interno del codice. Per la loro definizione si è optato per utilizzare delle strutture, le quali permetteranno di sfruttare una notazione compatta per le operazioni.

La prima struttura che si definisce è \texttt{tipo\_elemento}:

\begin{lstlisting}
module strutture
    implicit none
    
    type tipo_elemento
        integer :: nnodi, tipo_ele ! tipo ele da gmsh
        integer, dimension (:), pointer :: nodi ! array di nodi associato all elemento
        integer :: entity
    end type tipo_elemento
\end{lstlisting}

Al suo interno si trovano il numero di nodi che costituiscono l'elemento e il tipo di elemento che si ottiene dai tag di gmsh (es. quadrilatero o triangolare). Si ha inoltre un array contenente i tag univoci dei nodi associati all'elemento stesso.

Si definisce poi la struttura \texttt{tipo\_nodo}:

\begin{lstlisting}
    type tipo_nodo
        real (4), dimension (2) :: x
        integer :: nnghbrs, neles
        integer, dimension (:), allocatable :: nghbr
        integer, dimension (:), allocatable :: ele
    end type tipo_nodo
\end{lstlisting}

Al suo interno si trovano le coordinate $x$ e $y$ come elementi del vettore reale \texttt{x}. \texttt{nghbr} rappresenta una lista contenente i tag dei nodi vicini e \texttt{ele} è un vettore che contiene i tag degli elementi che si affacciano su quel nodo.

La terza struttura è \texttt{tipo\_interfaccia}:

\begin{lstlisting}
    type tipo_interfaccia
        integer :: e1, e2
        integer :: edge_e1, edge_e2
        integer :: nodo1, nodo2
        real (4), dimension (2) :: normal, x0
        real (4) :: length
        integer :: entity
        real (4), dimension (4) :: f
        integer :: int_perio ! 10 interno, sub bordi condizioni periodiche
    end type tipo_interfaccia
\end{lstlisting}

Al suo interno troviamo i tag degli elementi associati a quella interfaccia con la convenzione che $e_1$ è quello di riferimento per la normale uscente. Si hanno informazioni sui lati corrispondenti e i tag dei nodi agli estremi. Sono presenti quantità geometriche come il vettore normale ($n_x, n_y$), le coordinate del baricentro $x_0$ e la lunghezza dell'interfaccia. Vengono definiti anche i flussi che attraversano l'interfaccia (4 elementi per Eulero 2D).

L'ultima struttura è \texttt{tipo\_elemento\_solido}:

\begin{lstlisting}
    type tipo_elemento_solido
        integer :: nnodi, tipo_ele, nlati
        integer, dimension(:), allocatable :: nodi
        integer, dimension (:,:), allocatable :: edge
        integer :: nnghbrs
        integer, dimension(:), allocatable :: nghbr
        integer, dimension (:), allocatable :: interfaces
        real (4) :: area
        real (4), dimension (2) :: x0
        real (4), dimension (4) :: ucons
        real (4) :: u, v, a, P, T, S
        real (4), dimension (4) :: d_dt
    end type tipo_elemento_solido
\end{lstlisting}

La grandezza \texttt{ucons} contiene le 4 variabili conservative:
\[
ucons = \begin{Bmatrix} \rho \cdot E \\ \rho \\ \rho \cdot u \\ \rho \cdot v \end{Bmatrix}
\]

La grandezza \texttt{d\_dt} contiene le derivate temporali delle grandezze conservative:
\[
\frac{d}{dt} = \begin{Bmatrix} \frac{\partial (\rho E)}{\partial t} \\ \frac{\partial \rho}{\partial t} \\ \frac{\partial (\rho u)}{\partial t} \\ \frac{\partial (\rho v)}{\partial t} \end{Bmatrix}
\]

\section{Variabili}

Si definisce un nuovo modulo che richiama \texttt{strutture}:

\begin{lstlisting}
Module Variabili
    use strutture
    implicit none
    save
    
    integer :: k, kinf, kout, kfinal
    real (4) :: gam, GA, GB, GC, GD, GE, GF, GG, GH, GI, GJ
    real (4) :: time, dt, CFL, periodo
    character (100) :: mesh_file
    real (4) :: ttotin, ptotin, alpha, pexit, machin
    integer :: nnodi, ninterf, nele, nele_interni, nele_bordi, nentity, ndim, neqs
    real (4), dimension (4) :: norm2_residuals, norminf_residuals
    
    type(tipo_nodo), dimension(:), allocatable :: nodo
    type(tipo_interfaccia), dimension(:), allocatable :: interf
    type(tipo_elemento), dimension(:), allocatable :: ele_all
    type(tipo_elemento_solido), dimension(:), allocatable :: ele
    type(tipo_bordo), dimension(:), allocatable :: ele_bordi
    type(tipo_entity), dimension (6) :: entita
End module
\end{lstlisting}

\section{Main}

Il cuore del codice è il main, che richiama le diverse subroutine:

\begin{lstlisting}
Program test2d
    Use Variabili
    implicit none
    
    ! Pulizia file precedenti
    call system("del *.plt") ! Su Windows
    ! call system("rm *.plt") ! Su Linux
    
    call input
    call leggi_gmsh
    call processa_elementi
    call init
    
    call write_tecplot_file(0, 0.)
    time = 0.
    
    do k = 1, kfinal
        call compute_dt
        call compute_fluxes
        call integ_time
        time = time + dt
        
        if (mod(k, kinf) == 0) then
            write(*,*) 'k, time, dt = ', k, time, dt
            call compute_norm_residuals
        end if
        
        if (mod(k, kout) == 0) then
            call write_tecplot_file(k, 0.)
        end if
    end do
End Program
\end{lstlisting}

\subsection{Subroutine input}

Gestisce la lettura dei file \texttt{input.txt}, \texttt{inlet.txt} e \texttt{outlet.txt}.
\begin{itemize}
    \item \texttt{input.txt}: Nome mesh, \texttt{KFINAL}, \texttt{KINF}, \texttt{KOUT}, \texttt{CFL}.
    \item \texttt{inlet.txt}: $T^\circ_{IN}$, $P^\circ_{IN}$, $\alpha$, $M_{IN}$.
    \item \texttt{outlet.txt}: Pressione statica di uscita.
\end{itemize}

\subsection{Subroutine leggi\_gmsh}

Apre il file \texttt{.msh} e costruisce la connettività, allocando nodi ed elementi. Distingue tra elementi interni (triangoli, quadrilateri) e di bordo.

\subsection{Subroutine processa\_elementi}

Calcola le proprietà geometriche:
\begin{itemize}
    \item \textbf{Baricentro}: media delle coordinate dei nodi.
    \item \textbf{Area}: calcolata tramite la regola di Erone.
    \item \textbf{Lunghezza e Normale}: per le interfacce, calcolate tramite Pitagora e perpendicolarità al versore tangente.
\end{itemize}

Per un triangolo di lati $a, b, c$ e semiperimetro $p = (a+b+c)/2$:
\[ A = \sqrt{p(p-a)(p-b)(p-c)} \]

\subsection{Subroutine init}

Inizializza il campo di moto. Le equazioni di Eulero 2D in forma conservativa sono:
\[ \frac{\partial \{u\}}{\partial t} + \frac{\partial \{f\}}{\partial x} + \frac{\partial \{g\}}{\partial y} = 0 \]
Con:
\[ u = \begin{Bmatrix} \rho E \\ \rho \\ \rho u \\ \rho v \end{Bmatrix}, \quad f = \begin{Bmatrix} u(P+\rho E) \\ \rho u \\ P+\rho u^2 \\ \rho uv \end{Bmatrix}, \quad g = \begin{Bmatrix} v(P+\rho E) \\ \rho v \\ \rho uv \\ P+\rho v^2 \end{Bmatrix} \]

Le variabili vengono adimensionalizzate rispetto a valori di riferimento ($P_{IN}, T_{IN}$).

\section{Subroutine compute\_dt}

Calcola il passo temporale minimo basato sulla condizione CFL:
\[ \Delta t = CFL \frac{h}{|q| + a} \]
Dove $h \approx \sqrt{Area}$ e $\phi_{max} = |q| + a$ è la velocità massima di propagazione.

\section{Subroutine compute\_fluxes}

Valuta i flussi alle interfacce utilizzando solutori come Lax-Friedrichs o Roe. Integra le derivate temporali.

\subsection{Subroutine integ\_time}

Aggiorna le grandezze conservative usando Eulero esplicito:
\[ (\rho u)^{k+1}_i = (\rho u)^k_i + \Delta t \sum_{nlati} F_{ij} \cdot \hat{n}_{ij} \frac{l_{ij}}{A_i} \]

\begin{lstlisting}
subroutine integ_time
    use variabili
    implicit none
    integer :: i
    
    do i = 1, nele_interni
        ele(i)%ucons(:) = ele(i)%ucons(:) + dt * ele(i)%d_dt(:)
        ! Ricalcolo grandezze primitive dalle conservative
        ele(i)%u = ele(i)%ucons(3) / ele(i)%ucons(2)
        ele(i)%v = ele(i)%ucons(4) / ele(i)%ucons(2)
        ! ... altri ricalcoli
    end do
end subroutine
\end{lstlisting}
